\section{Graph Topology}
\label{section:22-graph-topology}
Graphs fit into core classes (scale-free, mesh, random) based on their size, shape, topological structure, and construction method. 
We begin by addressing the classes of graphs that occur in the real world, and then discuss synthetically generated graphs, which can be constructed randomly or using a prior distribution.

\subsection{Real-World Graphs}
This section will first identify the `real world' graph classes that are derived from real-world phenomena like email records, chat logs, network logs, road maps, protein structures, or any of a large collection of problems that can be represented as a graph.

\subsubsection{Scale-Free Graphs}
Scale-free graphs are a naturally occurring graph class characterized by a distribution of vertex degrees that follow a power-law distribution. 
Scale-free graphs are common in social networks because as the number of vertices increases over time, new vertices preferentially attach to vertices that already have high degrees~\cite{Barabasi1999a}. 
They tend to have a small diameter, a large variance in neighborhood size, and a `global-neighborhood' characteristic where within a few hops all vertices are in the same neighborhood.
These characteristics are often amalgamated together and termed the `small world' property.

\subsubsection{Mesh Graphs}
Mesh graphs are a naturally occurring graph class characterized by a flat, regular degree distribution.
Mesh graphs are often derived from physical phenomena like road networks, and so new vertices tend to extend the graph outwards~\cite{Beamer2015}.
They tend to have a large diameter, a small variance in neighborhood size and a `local-neighborhood' characteristic where only a small subset of vertices are reachable within a few hops of any given vertex.

\subsection{Synthetic Graphs}
Real-world graphs offer realistic test inputs for graph algorithms.
However, some domains lack sufficient data to support robust testing, especially regarding performance on irregular distributions. 
As a solution, synthetic graphs can be mathematically or imitatively generated to provide simple graphs for testing and evaluation. 

\subsubsection{Mathematically Generated Synthetic Graphs (Random Graphs)}
Multiple mathematical functions can be used to generate synthetic graphs, typically by adding vertices and edges following some set of heuristics to reflect a given probability distribution.
Notable generators include the Erd\H{o}s-R\'enyi generator to create \textit{Uniform Random Graphs}~\cite{Erdhos1960}, the \textit{Random Matrix} (RMAT)~\cite{Chakrabarti2004} and \textit{Kronecker}~\cite{Leskovec2010} generator to generate scale-free graphs.
Mathematically generated random graphs can be used to approximate real-world graphs but provide a sub-optimal substitute for the irregularities of real graphs because of the structural simplifications that occur in the generation process. 

\subsubsection{Imitatively Generated Synthetic Graphs (Imitative Graphs)}
Many graph datasets are not releasable for research and development because of privacy and confidentiality reasons. 
Active research in privacy preservation and synthetic graph generation seeks to statistically summarize a real-world graph, and then use that summary as the input to a synthetic graph generator, so that the synthetic graph has will have the same topological properties as the original, without the confidential information. 
Imitative graph generation methods use input features like \textit{Triangle Histograms, K-Hop Neighbor Histograms, Random Walk Histograms} and \textit{Number of Edges} to produce topologically similar graphs~\cite{Zahirnia2024}.
 process is inherently lossy and has only been tested on very small graphs (tens of thousands of edges), so is not yet reliable enough for wide use.


\subsection{Summary}

The core graph topology classes discussed above (scale-free, mesh, mathematically generated, imitatively generated) each have inherent benefits and drawbacks for graph processing. 
We view these classes as representatives of the kinds of `datasets' that are relevant to \gai. 
As such, a good \gai benchmark should cover each of these classes to enable a longitudinal view of expected performance.